\documentclass[12pt]{article}
\usepackage[a4paper,margin=2cm]{geometry}
\usepackage{amsmath,amssymb}
\usepackage{xepersian}
\settextfont[
  Path=../1401-2/HW2 - official solution/,
  Extension=.ttf,
  BoldFont=*Bd,
  ItalicFont=*It,
  BoldItalicFont=*BdIt
]{XB Niloofar}
\setdigitfont[
  Path=../1401-2/HW2 - official solution/,
  Extension=.ttf,
  BoldFont=*Bd,
  ItalicFont=*It,
  BoldItalicFont=*BdIt
]{XB Niloofar}
\setlength{\parindent}{0pt}
\setlength{\parskip}{0.55em}

\begin{document}
\begin{center}
  {\Large\textbf{اصلاحیهٔ پاسخ رسمی تمرین دوم ـ نیم‌سال ۱۳۹۹--۲}}
\end{center}

این برگه فقط موارد نادرست پاسخ منتشرشده را اصلاح می‌کند و باید همراه فایل اصلی خوانده شود.

\section*{سؤال ۱، بخش ب}
طبق صورت سؤال، محاسبات دو روش باید هم‌زمان و هنگامی متوقف شوند که شرط
\[
  |x_n-1|<0.02
\]
در \textbf{حداقل یکی} از دو روش برقرار شود؛ بنابراین تعداد تکرار دو روش باید یکسان باشد. پس از پنج تکرار داریم
\[
  x_{\mathrm{bisect}}=1.015625,
  \qquad
  x_{\mathrm{false\ position}}\simeq0.40787792.
\]
در این مرحله شرط توقف در روش تنصیف برقرار است و هر دو جدول باید متوقف شوند. ادامه‌دادن جدول روش نابجایی تا همگرایی مستقل آن، با خواستهٔ این بخش سازگار نیست.

اگر روش نابجایی را صرفاً برای بررسی جداگانه ادامه دهیم، نخستین تقریب تولیدشده‌ای که شرط را ارضا می‌کند، تقریب بیست‌وچهارم است:
\[
  x_{24}\simeq0.98350995.
\]
در جدول منتشرشده، مقدار تقریب چهاردهم باید
\(x_{14}=0.83787389\)
با
\(g(x_{14})=-0.82947558\)
باشد، نه
\(-0.88443\)؛
در نتیجه کران چپ سطر بعد نیز باید
\(a_{15}=0.83787389\)
شود. افزون بر این، سطر
\(x_{20}\)
از جدول حذف شده و دنباله‌ی درست انتهای جدول چنین است:
\[
\begin{array}{c|cc}
n&x_n&g(x_n)\\ \hline
20&0.95533398&-0.36678349\\
21&0.96494557&-0.30011260\\
22&0.97262969&-0.24233850\\
23&0.97871913&-0.19354380\\
24&0.98350995&-0.15318686
\end{array}
\]
بنابراین سطرهای شماره‌گذاری‌شده‌ی ۲۰ تا ۲۳ در نسخه‌ی منتشرشده در واقع شامل
\(x_{21}\)
تا
\(x_{24}\)
هستند و ستون‌های کران چپ و مقدار تابع نیز باید متناسب با آن جابه‌جا شوند.

\section*{سؤال ۱، بخش ج}
مرتبهٔ همگرایی روش نابجایی در حالت کلی \textbf{خطی} است. مرتبهٔ
\[
  \varphi=\frac{1+\sqrt5}{2}
\]
به روش \textbf{وتری} مربوط است، نه روش نابجایی. مقایسهٔ درست برای تابع
\(g(x)=x^{10}-1\)
پس از پنج تکرار نشان می‌دهد تنصیف شرط توقف را برآورده کرده است، در حالی که نابجایی به علت نامتوازن‌بودن مقادیر تابع در دو سر بازه بسیار کند پیش می‌رود.
\end{document}
